\section{相关工作}
\label{sec:related}

\nospacestitle{无自动扩缩容条件下的模型服务优化。 }
大规模模型服务并不简单。
大量研究关注如何高效利用 GPU 来加速模型服务~\cite{DBLP:conf/osdi/ZhongLCHZL0024,DBLP:journals/corr/abs-2404-09526,vllm,alpaserve,DBLP:conf/osdi/YuJKKC22,DBLP:conf/isca/PatelCZSGMB24,298501,kvcache}。
例如，Orca~\cite{DBLP:conf/osdi/YuJKKC22} 提出了迭代式调度与选择性批处理；
AlphaServe~\cite{alpaserve} 采用流水线并行来更好地应对负载尖峰，
但它无法动态调整流水线实例。
这些系统都假设运行在固定规模的 GPU 池上，
而我们已经说明了动态调整池规模的必要性，
并展示了如何通过 {\sys} 高效实现这一点。
{\sys} 为这些单实例模型服务系统补充了快速自动扩缩容机制：
我们建立在它们的高效单实例服务能力之上，
并在系统需要改变服务实例数量时提供超快扩缩容。

\stitle{动态扩缩服务实例。 }
动态扩缩服务实例的难点主要在于模型权重体积巨大且仍在持续增长，
因此把它们加载到加速器上（数据平面）十分耗时。
已有一些工作尝试加速该加载过程~\cite{pipeswitch, DBLP:conf/eurosys/JeongBA23, DBLP:conf/asplos/MiaoSDXL0J24, DBLP:conf/osdi/SunHZXZL024}。
例如，PipeSwitch~\cite{pipeswitch} 和 DeepPlan~\cite{DBLP:conf/eurosys/JeongBA23}
都利用模型逐层执行的特性，把推理与参数加载重叠起来以隐藏加载开销。
但它们只关注主机到设备的加载，而且当模型较大时，这种重叠仍然不是在线式的。
SpotServe~\cite{DBLP:conf/asplos/MiaoSDXL0J24} 与 Llumnix~\cite{DBLP:conf/osdi/SunHZXZL024}
实现了在线迁移，但迁移无法完全释放新旧两个实例的计算能力。
{\sys} 提供了一种新的机制，使服务实例在参数加载过程中也能在线扩缩，
即使加载尚未完成也能提升吞吐，
而我们已证明这对降低突发负载下的延迟至关重要。

并发工作 $\lambda$Scale~\cite{yu2025lambdascale} 也关注利用网络加速模型自动扩缩容。
关键区别在于，$\lambda$Scale 会把一次扩缩涉及的参数分散到多个实例上，
以降低扩缩时间，但代价是服务吞吐下降；
而 {\sys} 则可以以相近速度把完整参数无缝扩到所有实例上，
并借助基于多播链的扩缩过程避免牺牲吞吐。
此外，{\sys} 在扩缩过程中由于在线扩缩而具备逐步上升的吞吐，
而 $\lambda$Scale 仍然属于 stop-the-world 方法。

\stitle{加速无服务器计算中的冷启动。 }
加速模型扩缩容与无服务器计算中的冷启动加速一脉相承~\cite{sock,sand,FAASM,DBLP:conf/asplos/DuYXZYQWC20,DBLP:conf/asplos/Ustiugov0KBG21,DBLP:conf/eurosys/SaxenaJSKA22,DBLP:conf/usenix/WangCTWYLDC21}，
后者关注如何启动容器等通用计算实例。
我们继承了这些工作的思路，例如用于加速容器启动时间，
但进一步针对模型服务这一领域，结合其特定知识，
设计了高效的基于网络的在线自动扩缩容机制。
